\begin{textsample}{NEG Dim 1 – human – Score -9.00 – rj/rj\_ef\_linguagens\_lp\_1\_p11.txt}  \label{ex:f1_neg_001}
Absolutamente integrada a todas as outras áreas, a Língua Portuguesa assume de \textbf{forma} mais abrangente do que nunca o conceito de interdisciplinaridade para a formação integral do estudante.
Ser letrado é \textbf{ter} a \textbf{capacidade} de escutar, falar, ler e escrever sobre todos os assuntos presentes em qualquer componente curricular e nas questões pessoais ou \textbf{coletivas}, sempre atendendo às adaptações necessárias de comunicação, utilizando inclusive tecnologias digitais da informação e comunicação ( TDIC ), letramento digital e midiático. \textbf{Participar} do mundo letrado é fundamental quando a educação é \textbf{prevista} para todos e dessa \textbf{forma} o estudante \textbf{precisa} saber escutar, falar, ler e escrever, compreendendo, interpretando, \textbf{produzindo} e observando os aspectos culturais de sua língua materna, suas variantes, seus conhecimentos linguísticos, gramaticais, os gêneros textuais e suas estruturas, mas sempre à luz da linguagem.
Desta \textbf{maneira}, os estudantes podem ampliar sua participação na vida escolar e nos \textbf{grupos} sociais aos quais pertencem e/ou convivem, \textbf{possibilitando} que \textbf{participem} de \textbf{maneira} cada vez mais autônoma, crítica e reflexiva para que possam contribuir de \textbf{maneira} significativa para possíveis mudanças em suas práticas sociais cotidianas.

% matched lemmas: capacidade, coletivo, forma, grupo, maneira, participar, possibilitar, precisar, prever, produzir, ter
\end{textsample}
